\documentclass[10pt]{article}

\usepackage[margin=1.5cm]{geometry}
\usepackage{multicol}
\usepackage{booktabs}
\usepackage{array}
\usepackage{xcolor}
\usepackage{listings}
\usepackage{hyperref}

\definecolor{codegray}{rgb}{0.95,0.95,0.95}

\lstset{
    backgroundcolor=\color{codegray},
    basicstyle=\ttfamily\footnotesize,
    frame=single,
    breaklines=true,
    columns=fullflexible
}

\title{\textbf{Formulario de Bash Scripting}}
\author{}
\date{}

\begin{document}
\maketitle

\begin{multicols}{2}

%------------------------------------------------
\section*{Estructura básica}

\begin{lstlisting}
#!/bin/bash
echo "Hola mundo"
\end{lstlisting}

\begin{itemize}
\item \texttt{\#!/bin/bash}: intérprete
\item \texttt{chmod +x script.sh}
\item \texttt{./script.sh}
\end{itemize}

%------------------------------------------------
\section*{Variables}

\begin{lstlisting}
x=10
nombre="Rodrigo"
echo $x
echo "$nombre"
\end{lstlisting}

\begin{center}
\begin{tabular}{ll}
\toprule
Variable & Significado \\
\midrule
\texttt{\$0} & Nombre del script \\
\texttt{\$1..\$9} & Argumentos \\
\texttt{\$@} & Todos los argumentos \\
\texttt{\$?} & Código de salida \\
\texttt{\$\$} & PID del script \\
\bottomrule
\end{tabular}
\end{center}

%------------------------------------------------
\section*{Tipos de datos}

\begin{center}
\begin{tabular}{ll}
\toprule
Tipo lógico & Ejemplo \\
\midrule
String & \texttt{s="hola"} \\
Entero & \texttt{i=10} \\
Array & \texttt{a=(1 2 3)} \\
Array asociativo & \texttt{declare -A m} \\
\bottomrule
\end{tabular}
\end{center}

%------------------------------------------------
\section*{Arrays}

\begin{lstlisting}
arr=(uno dos tres)
echo ${arr[0]}
echo ${arr[@]}
\end{lstlisting}

\begin{lstlisting}
declare -A edades
edades[Juan]=20
\end{lstlisting}

%------------------------------------------------
\section*{Aritmética}

\begin{lstlisting}
a=10
b=3
echo $((a + b))
echo $((a * b))
\end{lstlisting}

%------------------------------------------------
\section*{Comparadores numéricos}

\begin{center}
\begin{tabular}{ll}
\toprule
Operador & Significado \\
\midrule
-eq & Igual \\
-ne & Distinto \\
-lt & Menor \\
-le & Menor o igual \\
-gt & Mayor \\
-ge & Mayor o igual \\
\bottomrule
\end{tabular}
\end{center}

\begin{lstlisting}
if [ $a -gt 5 ]; then
    echo "Mayor que 5"
fi
\end{lstlisting}

%------------------------------------------------
\section*{Comparadores de strings}

\begin{center}
\begin{tabular}{ll}
\toprule
Operador & Significado \\
\midrule
= & Igual \\
!= & Distinto \\
-z & Vacío \\
-n & No vacío \\
\bottomrule
\end{tabular}
\end{center}

%------------------------------------------------
\section*{Control de flujo}

\begin{lstlisting}
if [ $x -eq 10 ]; then
    echo "Diez"
elif [ $x -eq 5 ]; then
    echo "Cinco"
else
    echo "Otro"
fi
\end{lstlisting}

\begin{lstlisting}
for ((i=0;i<5;i++)); do
    echo $i
done
\end{lstlisting}

%------------------------------------------------
\section*{Funciones}

\begin{lstlisting}
suma() {
    echo $(($1 + $2))
}
resultado=$(suma 3 4)
\end{lstlisting}

%------------------------------------------------
\section*{Archivos}

\begin{center}
\begin{tabular}{ll}
\toprule
Prueba & Significado \\
\midrule
-e & Existe \\
-f & Archivo regular \\
-d & Directorio \\
-r & Lectura \\
-w & Escritura \\
-x & Ejecutable \\
\bottomrule
\end{tabular}
\end{center}

%------------------------------------------------
\section*{Redirecciones}

\begin{center}
\begin{tabular}{ll}
\toprule
Símbolo & Uso \\
\midrule
\texttt{>} & Sobrescribe \\
\texttt{>>} & Agrega \\
\texttt{|} & Pipe \\
\bottomrule
\end{tabular}
\end{center}

\begin{lstlisting}
ls | grep cpp > archivos.txt
\end{lstlisting}

%------------------------------------------------
\section*{Debug}

\begin{lstlisting}
bash -x script.sh
set -e
set -u
\end{lstlisting}

\end{multicols}

\end{document}
